\subsection{收敛级数}

4.1.1. 设 \(f = \sum_{\alpha} f_{\alpha}\) 是属于 \(\hat{\mathbf{P}}(\mathbf{E}_1, \ldots, \mathbf{E}_n; \mathbf{F})\) 的形式级数（参见附录）。若 \(\gamma\) 是 \(\mathbf{F}\) 上的连续半范数，且 \(\mathbf{R} = (\mathbf{R}_i)\) 是由 \(n\) 个 \(>0\) 的实数组成的系统，置

\[
\| f \| _ {\gamma , \mathbf {R}} = \sup _ {\alpha} \| f _ {\alpha} \| _ {\gamma} \mathbf {R} ^ {\alpha}.
\]

第 3.1.1 号（第二段）和第 3.1.2 号的定义与结果不加改变地适用；特别地，定义空间

\[
\mathcal {H} _ {\mathrm{R}} (\mathrm{E} _ {1}, \dots , \mathrm{E} _ {n}; \mathrm{F}) \quad \text { and } \quad \mathcal {H} (\mathrm{E} _ {1}, \dots , \mathrm{E} _ {n}; \mathrm{F}).
\]

4.1.2. \(j\) 到 \(\hat{\mathbf{P}}(\mathbf{E};\mathbf{F})\) 的规范同构 \(\hat{\mathbf{P}}(\mathbf{E}_{1},\ldots,\mathbf{E}_{n};\mathbf{F})\) 通过限制，对每个 \(\mathcal{H}_{\mathbb{R}}(\mathbf{E};\mathbf{F})\)，给出从 \(\mathcal{H}_{(\mathbb{R},\ldots,\mathbb{R})}(\mathbf{E}_{1},\ldots,\mathbf{E}_{n};\mathbf{F})\) 到 \(\mathbf{R}\in\mathbb{R}_{+}^{*}\) 的拓扑向量空间同构；它还给出从 \(\mathcal{H}(\mathbf{E};\mathbf{F})\) 到 \(\mathcal{H}(\mathbf{E}_{1},\ldots,\mathbf{E}_{n};\mathbf{F})\) 的同构。更确切地说，若 \(f=\sum_{m}f_{m}\in\hat{\mathbf{P}}(\mathbf{E};\mathbf{F})\) 且 \(j(f)=\sum_{\alpha}f_{\alpha}\)，则对 \(\gamma\) 上的每个连续半范数 \(\mathbf{F}\)，有：

\[
\begin{array}{l} \| f _ {m} \| _ {\gamma} = \sup _ {| \alpha | = m} \| f _ {\alpha} \| _ {\gamma} \\ \| f \| _ {\gamma , \mathbf {R}} = \| j (f) \| _ {\gamma , (\mathbf {R}, \dots , \mathbf {R})}. \end{array}
\]

4.1.3. 设 \(f = \sum_{\alpha} f_{\alpha}\) 是 \(\mathcal{H}(E_1, \ldots, E_n; F)\) 中的元素；令 \(I(f)\) 为满足如下条件的 \(R \in (\mathbf{R}_+^*)^n\) 的集合：对 F 上每个连续半范数 \(\gamma\)，当 \(\|f_{\alpha}\|_{\gamma} R^{\alpha}\) 趋于无穷时，乘积 \(|\alpha|\) 趋于零。集合 \(I(f)\) 非空；称其为 \(f\) 的严格收敛指标。点集 \(\Omega(f)\)

\[
(\log \mathbf {R} _ {1}, \dots , \log \mathbf {R} _ {n}) \quad \text { for } \quad \mathbf {R} \in I (f)
\]

是 \(\mathbf{R}^{n}\) 的凸子集。

当 \(n=1\) 时，集合 \(\mathbf{I}(f)\) 是 \(\mathbf{R}\) 中左开且右端开或闭的区间；其上界（有限或 \(+\infty\)）记作 \(\rho(f)\)，称为 \(f\) 的严格收敛半径。

使用 4.1.2 的记号，有 \(\mathbf{R} \in \mathbf{I}(f)\) 当且仅当

\[
(\mathbf {R}, \dots , \mathbf {R}) \in \mathrm{I} (j (f)).
\]

满足存在 \(x = (x_i)\) 且对 \(\mathbf{R} = (\mathbf{R}_i) \in I(f)\) 有 \(\| x_i \| \leqslant \mathbf{R}_i\) 的点 \(1 \leqslant i \leqslant n\) 的集合称为 \(f\) 的严格收敛域，记作 \(C(f)\)。它是 \(E\) 的开子集，是多球

\[
\mathrm{B} (\mathbf {R}) = \{x \in \mathrm{E} | \| x _ {i} \| \leqslant \mathrm{R} _ {i} \text {   for   } 1 \leqslant i \leqslant n \},
\]

当 \(\mathbf{R}\in\mathbf{I}(f)\) 时的并。

4.1.4. 将其中所有的 \(\tilde{\mathbf{C}}(f)\) 换为 \(\mathbf{C}(f)\)，将 \(\tilde{\mathcal{H}}_{\mathbb{R}}\) 换为 \(\mathcal{H}_{\mathbb{R}}\) 后，3.1.7 和 3.1.8 的结果仍然成立。

4.1.5. 设 \(\mathbf{F}_1, \ldots, \mathbf{F}_m\) 是完备赋范空间，且 \(\mathbf{F}\) 拟完备。令 \(\mathbf{f} = (f_i)_{1 \leqslant i \leqslant m}\)，其中 \(f_i \in \mathcal{H}(\mathbf{E}_1, \ldots, \mathbf{E}_n; \mathbf{F}_i)\)，并令 \(g \in \mathcal{H}(\mathbf{F}_1, \ldots, \mathbf{F}_m; \mathbf{F})\)，假设 \((f_i(0))_{1 \leqslant i \leqslant m}\) 中的点 \(\mathbf{E}\) 属于 \(g\) 的严格收敛域。于是，对每个 \(\alpha \in \mathbb{N}^m\)，形式级数 \(g_\alpha \circ \mathbf{f}\) 属于 \(\mathcal{H}(\mathbf{E}_1, \ldots, \mathbf{E}_n; \mathbf{F})\)，且族 \(g_\alpha \circ \mathbf{f}\) 在 \(\mathcal{H}(\mathbf{E}_1, \ldots, \mathbf{E}_n; \mathbf{F})\) 中可求和（从而当然也在 \(\hat{\mathbf{P}}(\mathbf{E}_1, \ldots, \mathbf{E}_n; \mathbf{F})\) 中可求和）。其和记作 \(g \circ \mathbf{f}\)。

更确切地说，存在 \(\mathbb{R} \in \bigcap_{i} \mathrm{I}(f_i)\) 和 \(\mathbb{R}' \in \mathrm{I}(g)\)，使得对 \(\sup_{|\alpha| > 0} \| f_{i,\alpha} \| \mathbb{R}^{\alpha} < \mathbb{R}_i'\) 有 \(1 \leqslant i \leqslant m\)。在这些条件下，形式级数 \(g_{\alpha} \circ \mathbf{f}\) 属于 \(\mathcal{H}_{\mathbb{R}}(\mathbf{E}_1, \ldots, \mathbf{E}_n; \mathbf{F})\)，且族 \(g_{\alpha} \circ \mathbf{f}\) 在 \(\mathcal{H}_{\mathbb{R}}(\mathbf{E}_1, \ldots, \mathbf{E}_n; \mathbf{F})\) 中可求和。最后，若 \(x \in \mathbf{B}(\mathbb{R})\)，则 \(\mathbf{f}(x) = (f_i(x))\) 属于 \(\mathbf{C}(g)\)，且有：

\[
g (\mathbf {f} (x)) = (g \circ \mathbf {f}) (x).
\]

再假设，对每个 \(i\)，存在 E＿i 中元素组成的族 \((e_{j}^{i})\)，使得 \(\mathbf{E}_{i}\) 中每个元素 \(x\)\(\mathbf{E}_{i}\) 都是可求和族 \((\lambda_{j}e_{j}^{i})\) 的和（其中 \(\lambda_{j}\in\mathbf{K}\)），且 \(\|x\|=\sup_{j}|\lambda_{j}|\)。于是，在前一段中可将条件 \(\sup_{|\alpha|>0}\|f_{i,\alpha}\|\mathbf{R}^{\alpha}<\mathbf{R}_{i}^{\prime}\) 换为条件 \(\|f_{i}\|_{\mathbf{R}}\leqslant\mathbf{R}_{i}^{\prime}\)。

4.1.6. 假设 \(E_{i} = K\)（\(1 \leqslant i \leqslant n\)）。于是将空间 \(\hat{\mathbf{P}}(\mathbf{K}^{n}; \mathbf{F})\) 识别为具有 F 中系数的 n 个不定元 \(X_{1}, \ldots, X_{n}\) 的形式幂级数空间，且 \(\hat{\mathbf{P}}(\mathbf{K}^{n}; \mathbf{F})\) 中的元素写作：

\[
f = \sum_ {\alpha} X ^ {\alpha} c _ {\alpha} \quad \text { with } c _ {\alpha} \in F.
\]

\[
\| f \| _ {\gamma , \mathbf {R}} = \sup \| c _ {\alpha} \| _ {\gamma}. \mathbf {R} ^ {\alpha}.
\]

若 \(\mathbf{R} \in (\mathbf{R}_{+}^{*})^{n}\)，且 \(\gamma\) 是 \(F\) 上的连续半范数，则有：

3.1.10 的最后一段仍然成立，3.1.11 也同样成立。

4.1.7. 假设 K 是完备赋值域（必为超度量的）L 的闭子域。对于 \(y \in E_{i} \otimes_{K} L\)，置：

\[
\| y \| = \inf (\sup _ {k} | a _ {k} |. \| x _ {k} \|)
\]

下确界取遍满足 \((x_{k}, a_{k}) \in E_{i} \times L\) 的所有有限对族 \(y = \sum_{k} x_{k} \otimes a_{k}\)。由此在 L-向量空间 \(E_{i} \otimes_{K} L\) 上得到一个范数，它在 \(E_{i}\) 上诱导出给定范数。用 \(E_{i}^{L}\) 表示 \(E_{i} \otimes_{K} L\) 关于此范数的完备化。

现在设 F 是分离完备的 L-向量多范数空间。对于 \(f_{\alpha}\) 上取值于作为 F 底层空间的 K-向量多范数空间 \(\alpha\) 的每个 K-多项式连续映射 \(E_1 \times \cdots \times E_n\)，若其关于多重次数 \(F_K\) 齐次，则存在唯一的在 \(\tilde{f}_{\alpha}\) 上取值于 F 的 L-多项式连续映射 \(E_1^L \times \cdots \times E_n^L\)，关于同一多重次数齐次，并延拓 \(f_{\alpha}\)。对 L-向量空间 F 上的每个连续半范数，有

\[
\| \tilde {f} _ {\alpha} \| _ {\gamma} = \| f _ {\alpha} \| _ {\gamma}.
\]

若 \(f = \sum_{\alpha} f_{\alpha} \in \mathcal{H}(E_1, \ldots, E_n; F_K)\)，则 \(\tilde{f} = \sum_{\alpha} \tilde{f}_{\alpha} \in \mathcal{H}(E_1^L, \ldots, E_n^L; F)\)。级数 \(f\) 和 \(\tilde{f}\) 具有相同的严格收敛指标（且当 \(n = 1\) 时具有相同的严格收敛半径）。

反之，令 L 为 K 的非离散闭子域，令 \(E_{i}^{0}\) 和 \(F^{0}\) 为分别从 \(E_{i}\) 和 F 限制标量得到的 L 上的空间。若 \(f = \sum f_{\alpha} \in \mathcal{H}(E_{1}, \ldots, E_{n}; F)\)，则 \(f_{\alpha} \in \mathrm{P}_{\alpha}(E_{1}^{0}, \ldots, E_{n}^{0}; F^{0})\)；若置 \(f^{0} = \sum_{\alpha} f_{\alpha} \in \hat{\mathrm{P}}(E_{1}^{0}, \ldots, E_{n}^{0}; F^{0})\)，则 \(f^{0} \in \mathcal{H}(E_{1}^{0}, \ldots, E_{n}^{0}; F^{0})\)。有 \(C(f) \subset C(f^{0})\)，且对所有 \(f(x) = f^{0}(x)\)，有 \(x \in C(f)\)。

